\documentclass{article}
\usepackage[preprint]{neurips_2026}
\usepackage[utf8]{inputenc}
\usepackage[T1]{fontenc}
\usepackage{hyperref}
\usepackage{url}
\usepackage{booktabs}
\usepackage{amsfonts}
\usepackage{nicefrac}
\usepackage{microtype}

\title{How Much of a VLA Leaderboard Score Is the Stack?}

\author{
  Yucheng Xie \\
  \texttt{xieyucheng@outlook.com}
}

\begin{document}

\maketitle

\begin{abstract}
VLA leaderboards compare models while rarely controlling what the scores
were produced \emph{with}: GPU type, weight dtype, benchmark image version,
seed set, sharding, asset mounting. We vary one knob at a time on a fixed
V100 stack and report each knob's score delta. Most knobs do not move scores
(benchmark image v0.7.0 vs v0.8.0: 86\% vs 92\%, n=50, within noise; rebuilt
slim image vs stock: 18/20 vs 18/20; external asset mount vs baked-in: 45/50
vs 46/50). One knob moves scores off a cliff: X-VLA in fp16 vs fp32 on
LIBERO goes from 0/30 to 29/30 across three suites --- the fp16 path fails
at the operator level (\texttt{float} vs \texttt{Half} bias mismatch on the
first inference), it does not merely degrade. We release a disclosure
checklist (GPU / dtype / image digest / seed protocol / shards / mounts) and
publish all images so every entry in our tables re-runs with one pull.
\end{abstract}

\section{Introduction}
A VLA score is a joint product of model, benchmark, and stack. The community
standardizes the first two and ignores the third. Papers and leaderboards
routinely compare an fp16 number from one lab against a bf16 number from
another, or a v0.7.0 number against a v0.8.0 number, as if the stack were
transparent. We test whether it is.

\section{Method}
One knob at a time, everything else pinned (V100 32\,GB, fixed seeds, fixed
shards unless the knob \emph{is} sharding). Artifacts (aggregate JSON,
SQLite recordings, logs) are archived per run; every image is published with
its Dockerfile. Negative-control discipline: mounts are verified by breaking
them on purpose; image identity is verified by digest, not tag.

\section{Results}
\label{sec:results}

\paragraph{Benchmark image (v0.7.0 vs v0.8.0).} grab\_roller $\times$50, same
seeds: 43/50 (86\%) vs 46/50 (92\%). Difference 6pp against a ${\sim}12$pp
2-sigma band at n=50: no significant difference. Scope: one task.

\paragraph{Baked vs native.} Rebuilt slim image (assets removed, nvcc
stripped, 11.3\,GB vs 21\,GB) vs the upstream-identical stock image,
grab\_roller $\times$20 same seeds: 18/20 vs 18/20. Zero delta; 50-episode
equivalence 46/50 vs 46/50 exact.

\paragraph{Asset mounting.} External 16\,GB asset mount vs image-baked
assets, same 50 seeds: 45/50 (90\%) vs 46/50 (92\%). Equivalent.

\paragraph{Precision (the cliff).} X-VLA on LIBERO (spatial/object/goal
$\times$10, same seed): fp32 scores 9/10, 10/10, 10/10 (29/30); fp16 scores
0/10 $\times$3 (0/30). The fp16 zeros are not low scores: every episode
errors at the first inference with
\texttt{RuntimeError: Input type (float) and bias type (c10::Half) should be
the same}. The denoising path is numerically broken in fp16 on this stack.
A 7\,B CogACT in fp32 does not load on a 32\,GB card (OOM at first forward);
sharded across two cards it loads, but its custom forward uses bare
cross-device concatenations, so split inference needs model surgery ---
reported as a limit, not a result.

\paragraph{Sharding.} 2 shards stable, 4 slow, 8 hangs. Speed/stability only.

\paragraph{Seeds.} Blind seeds vs bundled verified seeds move the error rate,
not the valid score; re-running the same bundled 50 seeds gives 41/50 vs
46/50 (diff $-10$pp, ${\sim}1.7\sigma$: noise).

\paragraph{Cross-stack reference.} place\_empty\_cup: our V100 fp16 16--20\%
vs official H100 bf16 28\%. Recorded as pending (different GPU \emph{and}
dtype), not judged.

\section{Related work}
vla-eval (protocol layer), VLATest (environment perturbations),
VLA-Perf (cross-GPU speed only, assumes BF16/FP16 interchangeable --- \S\ref{sec:results}
falsifies that for at least X-VLA), leaderboard-design studies, REPLICA
(real-robot cross-site), practical eval guides (checklists, no precision
study), ROCm accuracy-parity work (hardware direction). None quantifies
same-code-different-stack score deltas. That is this paper.

\section{The disclosure checklist}
GPU model + driver; weight dtype; benchmark image digest (+ slim/patch
notes); serve image digest; seed protocol; shard count; asset/code mounts;
harness ref. Missing one item: not comparable.

\section{Conclusion}
Stacks are not transparent: one common knob (fp16) takes a model from 97\%
to 0\% without changing a single weight. Compare same-stack, or disclose.

\section*{Acknowledgments}
Compute provided by Huawei Cloud. AI coding assistance was used for
engineering (experiment orchestration, plotting) under human direction; all
results were produced by executed runs.

\end{document}
